%=======================================================================
%  CSE558 Data Science - Midterm Project Report
%  Monsoon 2026, IIIT-Delhi
%
%  Built on the course template. Text in LIGHT BLUE (\eg{...}) is still a
%  placeholder: team names, roll numbers, emails and the two links.
%  Every number comes from results/*.json, produced by code/run_all.py.
%  Figures are drawn by pgfplots from figdata/*.csv.
%  Compile with pdfLaTeX, twice.
%=======================================================================
\documentclass[10pt,a4paper,twocolumn]{article}

%========================= EDIT THIS =========================
% Your project title: it sets the heading AND the PDF's title.
\newcommand{\projecttitle}{Lit roads, lighter crashes? Street lighting,
  motorcycles and crash severity in UK accident records}
\newif\ifshowhints
\showhintsfalse     % \showhintsfalse hides every grey instruction line
%=============================================================

%-------------------------------------------- fonts (no need to edit)
\usepackage{iftex}
\ifPDFTeX
  \usepackage[T1]{fontenc}
  \usepackage[utf8]{inputenc}
  \IfFileExists{carlito.sty}%
    {\usepackage[sfdefault]{carlito}}%              Calibri look
    {\usepackage[scaled=0.92]{helvet}\renewcommand{\familydefault}{\sfdefault}}%
  \usepackage{courier}
\else
  \usepackage[no-math]{fontspec}
  \setmainfont{Carlito}[Extension=.ttf, UprightFont=*-Regular, BoldFont=*-Bold,
    ItalicFont=*-Italic, BoldItalicFont=*-BoldItalic]
\fi

%----------------------------------------------------------- packages
\usepackage[margin=18mm,bottom=20mm,columnsep=7mm]{geometry}
\usepackage{amsmath}
\usepackage{xcolor}
\usepackage{graphicx}
\usepackage{tabularx}
\usepackage{array}
\usepackage{booktabs}
\usepackage{enumitem}
\usepackage{titlesec}
\usepackage{needspace}
\usepackage[font=small,labelfont=bf]{caption}
\usepackage{pgfplots}
\pgfplotsset{compat=1.18, table/col sep=comma}
\usepackage[hidelinks]{hyperref}
\hypersetup{pdftitle={\projecttitle},
            pdfsubject={CSE558 Data Science, Monsoon 2026: Midterm Project Report}}

%------------------------------------------------------------- layout
\definecolor{eg}{HTML}{7396C4}         % worked example (light blue)
\definecolor{hintgrey}{gray}{0.58}     % instructions (light grey)
\definecolor{cA}{HTML}{2F5D8C}         % plot colour 1
\definecolor{cB}{HTML}{D9822B}         % plot colour 2

\titleformat{\section}{\normalfont\large\bfseries}{\thesection.}{0.5em}{}
\titlespacing*{\section}{0pt}{12pt}{3pt}
\titleformat{\subsection}{\normalfont\normalsize\bfseries}{\thesubsection}{0.5em}{}
\titlespacing*{\subsection}{0pt}{10pt}{2pt}
\newcommand{\sectionbreak}{\needspace{10\baselineskip}}
\newcommand{\subsectionbreak}{\needspace{10\baselineskip}}

\setlength{\parindent}{0pt}
\setlength{\parskip}{4pt}
\setlist{leftmargin=1.2em, itemsep=2pt, topsep=2pt, parsep=0pt}
\renewcommand{\arraystretch}{1.25}
\setlength{\tabcolsep}{4pt}
\newcolumntype{L}[1]{>{\raggedright\arraybackslash}p{#1}}
\newcolumntype{Y}{>{\raggedright\arraybackslash}X}
\newcolumntype{K}[1]{>{\raggedright\arraybackslash\bfseries\boldmath}p{#1}}  % bold labels

%--------------------------------------------------------- shorthands
\newcommand{\hint}[1]{\ifshowhints{\footnotesize\itshape\color{hintgrey}#1\par}\fi}
\newcommand{\eg}[1]{\textcolor{eg}{#1}}
\newcommand{\egblock}[1]{{\color{eg}#1\par}}
\newcommand{\field}[2]{\textbf{#1}~#2\par}
\newcommand{\tpart}[2]{\par\needspace{12\baselineskip}\smallskip\textbf{#1)~#2}\par}
\newcommand{\minihead}[1]{\par\needspace{8\baselineskip}\smallskip\textbf{#1}\par}
\newcommand{\brkus}{\textunderscore\allowbreak}
\newcommand{\code}[1]{{\ttfamily\let\_\brkus #1}}   % may break after _
\newcommand{\pp}{\,pp}
\newcommand{\figfile}[2]{%
  \IfFileExists{#1}%
    {\includegraphics[width=\linewidth,height=#2,keepaspectratio]{#1}}%
    {\fbox{\parbox[c][#2][c]{\dimexpr\linewidth-2\fboxsep-2\fboxrule\relax}%
      {\centering\small\itshape\color{hintgrey}Your plot: \detokenize{#1}}}}}

% shared look for every plot
\pgfplotsset{every axis/.append style={
  width=\linewidth, font=\footnotesize, axis lines*=left,
  tick style={draw=none}, ymajorgrids, grid style={gray!25},
  legend style={draw=none, font=\footnotesize, fill=none},
  every axis label/.append style={font=\footnotesize}}}

%=======================================================================
\begin{document}

\twocolumn[%
{\small\bfseries CSE558 Data Science\ \ \textbullet\ \ Midterm Project Report\ \ \textbullet\ \ Monsoon 2026\par}
\vspace{4pt}
{\LARGE\bfseries\raggedright\projecttitle\par}
\vspace{6pt}
\rule{\linewidth}{1pt}

{\small
\begin{tabularx}{\linewidth}{@{}K{2.4cm} Y@{}}
\toprule
Group name & \eg{Team name} \\
Dataset & UK road accidents, \code{accident\_data.csv} (Kaggle upload of DfT STATS19 police records): \eg{kaggle.com/datasets/\dots} \\
Code & \eg{github.com/\dots/cse558-midterm} (Python scripts in \code{code/}, \code{README.md}, every reported number in \code{results/}) \\
\bottomrule
\end{tabularx}}

\vspace{2pt}
{\small
\begin{tabularx}{\linewidth}{@{}L{0.5cm} L{3.4cm} L{2.3cm} L{4.8cm} Y@{}}
\toprule
\textbf{\#} & \textbf{Name} & \textbf{Roll no.} & \textbf{Institute email} & \textbf{Contribution} \\
\midrule
1 & \eg{Name} & \eg{Roll no.} & \eg{email@iiitd.ac.in} & \eg{Dataset, preparation (1--3)} \\
2 & \eg{Name} & \eg{Roll no.} & \eg{email@iiitd.ac.in} & \eg{EDA, problem statements (4)} \\
3 & \eg{Name} & \eg{Roll no.} & \eg{email@iiitd.ac.in} & \eg{Tests 1--2 (5)} \\
4 & \eg{Name} & \eg{Roll no.} & \eg{email@iiitd.ac.in} & \eg{Tests 3--4, limitations, code (5, 6)} \\
\bottomrule
\end{tabularx}}
\vspace{12pt}
]

%-----------------------------------------------------------------------
\section{Dataset}

\field{Source:}{a Kaggle upload of the UK Department for Transport's STATS19 records of
police-reported road accidents with at least one injury in Great Britain. One row is one
accident. The accident IDs (e.g.\ \code{200701BS64157}) are STATS19 accident indices, with the
year in the first four digits.}

\field{What is raw about it:}{the accident ID was turned into scientific notation
(\code{2.01E+07}) for 239{,}884 rows, which creates 239{,}659 false duplicate IDs; for every
readable 2007--2010 ID the date is exactly 12 years later than the ID year; six district names lost the letters ``na''
(\code{Stevege} for Stevenage); 19{,}739 cells are missing; and Fatal is only 1.3\% of rows.}

{\small
\begin{tabularx}{\linewidth}{@{}L{3.0cm} Y Y@{}}
\toprule
 & \textbf{Raw file} & \textbf{After encoding} \\
\midrule
Instances (rows) & 670{,}679 & 670{,}660 \\
Attributes (columns) & 14 & 52 \\
Instances $\times$ attributes & 9{,}389{,}506 & 34{,}874{,}320 \\
\bottomrule
\end{tabularx}}

\field{Size check:}{$670{,}660\times52=34{,}874{,}320\ge3{,}000{,}000$, so the requirement is
met (the raw file alone gives $670{,}679\times14=9{,}389{,}506$). The 52 columns are 4 numeric,
severity as an ordinal, the KSI flag, district frequency, and 45 one-hot columns: light 5, road
surface 6, road type 6, urban/rural 3, weather 9, vehicle type 16.}

\minihead{Key features}
{\footnotesize
\begin{tabularx}{\linewidth}{@{}L{2.5cm} L{1.4cm} Y@{}}
\toprule
\textbf{Feature} & \textbf{Type} & \textbf{Real-world meaning} \\
\midrule
\code{Accident\_Severity} & ordinal & worst injury: Slight, Serious or Fatal. Our response is
\textbf{KSI} = Fatal or Serious (14.7\% of crashes), the standard UK road-safety measure \\
\code{Light\_Conditions} & nominal & daylight, or darkness with lights lit, unlit, none or
unknown \\
\code{Urban\_or\_Rural\_Area} & nominal & urban or rural location \\
\code{Road\_Type} & nominal & single or dual carriageway, roundabout, one-way, slip road \\
\code{Road\_Surface\_Conditions} & nominal & dry, wet, frost/ice, snow, flood \\
\code{Vehicle\_Type} & nominal & vehicle recorded for the crash, 15 types \\
\code{Number\_of\_Casualties} & discrete & people injured, 1 to 68 \\
\code{Number\_of\_Vehicles} & discrete & vehicles involved, 1 to 32 \\
\code{District Area} & nominal & local-authority district, 416 levels after cleaning \\
\bottomrule
\end{tabularx}}

%-----------------------------------------------------------------------
\section{Data preparation}

{\footnotesize
\begin{tabularx}{\linewidth}{@{}L{1.9cm} L{2.5cm} Y@{}}
\toprule
\textbf{Issue found} & \textbf{What we did} & \textbf{Why, and not the alternative} \\
\midrule
19 exact duplicate rows
 & dropped
 & the same crash counted twice would add weight to it in every test \\
ID in scientific notation for 239{,}884 rows
 & kept the ID out of every key and every test
 & de-duplicating on the ID would have deleted 239{,}659 distinct crashes \\
dates 12 years ahead of the ID year (all 420{,}776 readable older IDs)
 & date not used, and dropped from the encoding
 & no test needs time; shifting back would require guessing for the IDs that are unreadable \\
6 district names missing ``na'' (3{,}438 rows)
 & merged back: \code{Stevege} $\to$ Stevenage, \code{Blaeu Gwent}, \code{North/South
 Larkshire}, \code{Clackmannshire}, \code{Crewe and ntwich}
 & Test 2 pairs crashes within districts; a split district halves its counts and adds a fake
 one \\
missing categories: weather 14{,}351, road type 4{,}580, surface 742, urban/rural 26, vehicle 7
 & recoded to an explicit \code{Unknown} level
 & filling with the mode would push 4{,}580 unknown road types into ``single carriageway'',
 which is PS1's filter; \code{Unknown} rows simply fall out of each filter \\
51 missing coordinates
 & kept as missing
 & coordinates are not used \\
casualty outliers: IQR fence is $Q_3+1.5\,\mathrm{IQR}=1$
 & kept all values (max 68)
 & a fence at 1 would flag all 158{,}994 multi-casualty crashes; they are real (e.g.\ bus
 crashes). Test 4 is re-run without the 67 crashes above 10 casualties \\
vehicles above fence 3.5 (15{,}310)
 & kept
 & real pile-ups; PS2 fixes \code{Number\_of\_Vehicles}$\,=2$ anyway \\
6 nominal columns
 & one-hot (45 columns); severity ordinal $0/1/2$ plus a KSI flag
 & categories have no order; KSI is the response \\
district, 416 levels
 & frequency-encoded
 & one-hot would add 415 sparse columns; the raw label is kept for Test 2 \\
\bottomrule
\end{tabularx}}

%-----------------------------------------------------------------------
\section{What changed after preparation}

{\footnotesize
\begin{tabularx}{\linewidth}{@{}L{2.9cm} L{1.7cm} Y@{}}
\toprule
\textbf{Measure} & \textbf{Before} & \textbf{After} \\
\midrule
Rows & 670{,}679 & 670{,}660 \\
Attributes & 14 & 52 (encoded) \\
Missing cells & 19{,}739 & 0 in the columns we use; 51 coordinates left, unused \\
Duplicate rows & 19 & 0 \\
Duplicate IDs & 239{,}659 & ID no longer used as a key \\
District levels & 422 & 416 (6 typo levels merged) \\
Usable dates & 0 (shifted) & dropped \\
KSI share & 14.67\% & 14.67\% (unchanged) \\
\bottomrule
\end{tabularx}}

Turning severity into the binary KSI flag is what makes the proportion tests possible. Merging
the district typos moved 3{,}438 crashes back to their real district, which changes the
district-level pairs in Test 2. Only 19 rows were lost, so no group we test was thinned out by
cleaning.

%-----------------------------------------------------------------------
\section{Basic questions and problem statement}

\begin{figure}[!ht]
\centering
\begin{tikzpicture}
\begin{axis}[ybar, height=3.9cm, bar width=7pt, ymin=0, ymax=25,
  ylabel={KSI share (\%)},
  xtick=data, xticklabels from table={figdata/ksi_by_light.csv}{light},
  enlarge x limits=0.13,
  legend style={at={(0.5,1.03)}, anchor=south}, legend columns=2]
\addplot[fill=cA, draw=none] table[x expr=\coordindex, y=Rural, col sep=comma]{figdata/ksi_by_light.csv};
\addplot[fill=cB, draw=none] table[x expr=\coordindex, y=Urban, col sep=comma]{figdata/ksi_by_light.csv};
\legend{Rural, Urban}
\end{axis}
\end{tikzpicture}
\caption{KSI share by light condition (all except Daylight are in darkness). Rural crashes on
unlit roads are the most severe group: 23.8\% KSI against 17.9\% where lights are lit.}
\label{fig:light}
\end{figure}

\begin{figure}[!ht]
\centering
\begin{tikzpicture}
\begin{axis}[xbar, width=0.82\linewidth, height=3.7cm, bar width=6pt, xmin=0, xmax=0.15,
  xlabel={Cram\'er's $V$ with KSI},
  ytick=data, yticklabels from table={figdata/cramers_v_ksi.csv}{column},
  enlarge y limits=0.1, xmajorgrids, ymajorgrids=false,
  scaled x ticks=false, xticklabel style={/pgf/number format/fixed},
  nodes near coords, nodes near coords style={font=\tiny, /pgf/number format/fixed,
    /pgf/number format/precision=3}]
\addplot[fill=cA, draw=none] table[y expr=\coordindex, x=v, col sep=comma]{figdata/cramers_v_ksi.csv};
\end{axis}
\end{tikzpicture}
\caption{How strongly each column is associated with KSI. Vehicle type is the weakest by far
($V=0.005$), 13 times weaker than light conditions. This warns us before PS2 is tested.}
\label{fig:cramer}
\end{figure}

\begin{figure}[!ht]
\centering
\begin{tikzpicture}
\begin{axis}[ybar, height=3.5cm, bar width=8pt, ymin=0, ymax=85,
  ylabel={\% of crashes}, xlabel={casualties per crash (two-vehicle crashes)},
  xtick=data,
  xticklabels from table={figdata/casualty_dist.csv}{casualties},
  enlarge x limits=0.15, legend pos=north east]
\addplot[fill=cA, draw=none] table[x expr=\coordindex, y=Motorcycle, col sep=comma]{figdata/casualty_dist.csv};
\addplot[fill=cB, draw=none] table[x expr=\coordindex, y=Car, col sep=comma]{figdata/casualty_dist.csv};
\legend{Motorcycle, Car}
\end{axis}
\end{tikzpicture}
\caption{Casualties per crash are almost identical for motorcycles and cars: about 74\% of
crashes hurt exactly one person, and the counts are heavily right-skewed (skewness 7.0 over all
crashes).}
\label{fig:cas}
\end{figure}

\minihead{Basic questions}
\begin{enumerate}[label=\arabic*.]
  \item Which vehicle type is in the most crashes? Cars and taxis, in 77.4\% of them.
  \item Are rural crashes more severe than urban ones? Yes: 18.0\% KSI against 12.8\%.
  \item Are crashes in the dark more severe than in daylight? Yes, in both areas
  (Figure~\ref{fig:light}).
  \item Does bad weather make crashes worse? No: snowfall (9.7\%) and icy roads (11.9\%) have
  \emph{lower} KSI shares than dry roads (15.1\%), probably because people drive slower.
\end{enumerate}

\minihead{Our problem statements}

\noindent\fbox{\begin{minipage}{\dimexpr\linewidth-2\fboxsep-2\fboxrule\relax}
\textbf{PS1 (primary).} Among rural single-carriageway crashes in darkness, is the KSI share
higher where there is no street lighting than where lights are lit? Councils pay to light rural
roads; if unlit crashes are no more severe, that spend is hard to justify on severity grounds.

\smallskip
\textbf{PS2.} In two-vehicle crashes, are motorcycle crashes more often KSI, while hurting
\emph{fewer} people per crash, than car crashes? If so, budgets split by casualty counts would
understate motorcycle safety.
\end{minipage}}

\field{How they are novel:}{the basic questions compare one column against severity. PS1
compares two \emph{darkness} groups on the same road type, so darkness itself is held fixed and
only the lighting differs; it then checks the result within districts. PS2 asks two questions
about the same groups, one on severity and one on casualty count, which could point in opposite
directions. A plain ``which vehicle has the most crashes'' cannot show that.}

%-----------------------------------------------------------------------
\section{Statistical inferences}

\field{Rows tested, PS1:}{170{,}583 rural single-carriageway crashes, of which 28{,}201 were in
darkness with no lighting ($N$) and 14{,}382 in darkness with lights lit ($L$). Excluded:
``lights unlit'' (608) and ``lighting unknown'' (1{,}683), whose lighting status is unclear,
and daylight.}

\field{Rows tested, PS2:}{397{,}950 two-vehicle crashes, of which 34{,}017 are motorcycles of
any engine size ($M$) and 307{,}722 are cars or taxis ($C$). Other vehicles are excluded.}

%------------------------------------------------------------ Test 1
\subsection*{Test 1 (PS1): is the KSI share higher on unlit rural roads at night?}

\tpart{a}{Hypotheses, written before running the test}
{\small
\begin{tabularx}{\linewidth}{@{}K{3.4cm} Y@{}}
\toprule
$H_0$ & $p_N-p_L\le0$ \\
$H_1$ & $p_N-p_L>0$ \\
Significance level & $\alpha=0.05$ \\
Side & one-sided (right) \\
Test & two-proportion $z$-test (pooled SE for $z$, unpooled SE for the CI) \\
Smallest difference that matters & 2 percentage points \\
\bottomrule
\end{tabularx}}

$p_N$, $p_L$: KSI share of crashes on unlit and lit roads. One-sided, because the claim is that
no lighting is worse. 2\pp\ means 20 more KSI crashes per 1{,}000 rural night crashes. That is roughly where the cost of a lighting scheme becomes worth discussing.

\tpart{b}{Other tests we considered}
Two independent groups and a binary outcome, with expected cell counts in the thousands, call
for a large-sample test on the difference of two proportions.

{\footnotesize
\begin{tabularx}{\linewidth}{@{}L{2.4cm} Y@{}}
\toprule
\textbf{Also considered} & \textbf{Why we did not use it} \\
\midrule
Chi-square, $2\times2$ & identical evidence ($\chi^2=65.32=z^2$, $p=6.4\times10^{-16}$
two-sided) but two-sided only, and it gives no CI for the difference \\
Fisher's exact test & needed only for small expected counts; our smallest is 3{,}348 \\
CMH test by road surface & valid and stricter, since wet nights could confound. Run as a
robustness check (below) \\
Logistic regression & adjusts for many confounders at once, but it is a model rather than one
of our tests; left as an extension \\
\bottomrule
\end{tabularx}}

\tpart{c}{Result}
{\small
\begin{tabularx}{\linewidth}{@{}K{3.4cm} Y@{}}
\toprule
Estimate & $+3.50$\pp\ (24.46\% unlit, 20.96\% lit) \\
95\% confidence interval & $+2.67$ to $+4.33$\pp \\
Test statistic (d.f.) & $z=8.08$ (normal; $\chi^2$ form has 1 d.f.) \\
p-value & $3.2\times10^{-16}$ \\
Decision & reject $H_0$ \\
\bottomrule
\end{tabularx}}

Rural night crashes on unlit roads are 3.5\pp\ more likely to be KSI, an odds ratio of 1.22.
\textbf{It matters:} the whole CI lies above the 2\pp\ threshold. \emph{Robustness:} stratifying
by road surface (dry, wet, frost/ice, snow, flood), the CMH test gives $z=9.61$,
$p=3.7\times10^{-22}$, with a Mantel--Haenszel risk difference of $+4.18$\pp, so road surface
does not explain the gap. The gap is concentrated on dry (28.7\% vs 23.6\%) and wet (23.4\% vs
19.3\%) roads, and it disappears on frost or ice (14.4\% vs 14.6\%), where everyone already drives
slowly.

%------------------------------------------------------------ Test 2
\subsection*{Test 2 (PS1): does the gap hold within districts?}

Test~1 treats every crash as independent, but crashes in one district share roads, policing and
reporting habits. Here each district is its own control: for every district with at least 20
crashes in \emph{both} groups we take $d_i=\text{KSI}_N-\text{KSI}_L$.

\tpart{a}{Hypotheses, written before running the test}
{\small
\begin{tabularx}{\linewidth}{@{}K{3.4cm} Y@{}}
\toprule
$H_0$ & $\mu_d\le0$ \\
$H_1$ & $\mu_d>0$ \\
Significance level & $\alpha=0.05$ \\
Side & one-sided (right) \\
Test & paired $t$-test on the district differences $d_i$ \\
Smallest difference that matters & 2 percentage points \\
\bottomrule
\end{tabularx}}

\tpart{b}{Other tests we considered}
Each district gives a pair of shares, so the test should use the within-district difference.
The $d_i$ look normal (skewness 0.10, Shapiro--Wilk $p=0.72$), so a $t$-test on them is
appropriate.

{\footnotesize
\begin{tabularx}{\linewidth}{@{}L{2.4cm} Y@{}}
\toprule
\textbf{Also considered} & \textbf{Why we did not use it} \\
\midrule
Wilcoxon signed-rank & meant for skewed $d_i$ or outliers, which we do not have. Run alongside:
$W=20{,}981$, $p=2.3\times10^{-7}$ \\
Sign test & valid, but uses only the sign of each $d_i$, so it has less power: 158 of 247
positive, $p=6.7\times10^{-6}$ \\
Unpaired Welch $t$-test & ignores the pairing and leaves district-to-district variation in the
error: $t=4.80$, $p=1.1\times10^{-6}$, weaker than the paired $t=5.57$ because the two shares
are correlated across districts ($r=0.26$) \\
\bottomrule
\end{tabularx}}

\tpart{c}{Result}
{\small
\begin{tabularx}{\linewidth}{@{}K{3.4cm} Y@{}}
\toprule
Estimate & mean $d=+3.08$\pp\ (median $+2.93$\pp; 247 of 396 districts, 37{,}189 crashes) \\
95\% confidence interval & $+1.99$ to $+4.16$\pp \\
Test statistic (d.f.) & $t=5.57$ (246) \\
p-value & $3.3\times10^{-8}$ \\
Decision & reject $H_0$ \\
\bottomrule
\end{tabularx}}

\begin{figure}[!ht]
\centering
\begin{tikzpicture}
\begin{axis}[ybar, height=3.4cm, bar width=4.4, ymin=0,
  xlabel={$d_i$ = KSI share unlit $-$ lit (pp)}, ylabel={districts},
  xmin=-25, xmax=30, xtick={-20,-10,0,10,20}]
\addplot[fill=cA, draw=none] table[x=mid, y=count, col sep=comma]{figdata/district_diff_hist.csv};
\draw[cB, very thick] (axis cs:3.08,0) -- (axis cs:3.08,60);
\end{axis}
\end{tikzpicture}
\caption{District differences $d_i$ (247 districts). They are spread widely (SD 8.7\pp), but
centred above zero; the orange line is the mean, $+3.1$\pp.}
\label{fig:dist}
\end{figure}

The within-district comparison gives the same answer as Test~1, so the result survives the
clustering criticism. \textbf{It probably matters:} the estimate of 3.1\pp\ is above the
threshold, but the lower end of the CI (1.99\pp) touches 2\pp, so a gap just below the threshold
cannot be ruled out at district level.

%------------------------------------------------------------ Test 3
\subsection*{Test 3 (PS2): are motorcycle crashes more often KSI?}

\tpart{a}{Hypotheses, written before running the test}
{\small
\begin{tabularx}{\linewidth}{@{}K{3.4cm} Y@{}}
\toprule
$H_0$ & $p_M-p_C\le0$ \\
$H_1$ & $p_M-p_C>0$ \\
Significance level & $\alpha=0.05$ \\
Side & one-sided (right) \\
Test & two-proportion $z$-test \\
Smallest difference that matters & 5 percentage points \\
\bottomrule
\end{tabularx}}

$p_M$, $p_C$: KSI share of two-vehicle crashes recorded with a motorcycle or a car. A gap below
5\pp\ would not justify a separate motorcycle-safety budget line.

\tpart{b}{Other tests we considered}
Again two independent groups and a binary outcome with large counts.

{\footnotesize
\begin{tabularx}{\linewidth}{@{}L{2.4cm} Y@{}}
\toprule
\textbf{Also considered} & \textbf{Why we did not use it} \\
\midrule
Chi-square, severity (3 levels) $\times$ group & uses all three levels, but the Fatal cells are
thin and the result is harder to read. Supporting run: $\chi^2=1.65$ (2 d.f.), $p=0.44$ \\
Fisher's exact test & counts are far too large to need it (smallest expected count 301) \\
Relative risk with CI & useful alongside, but our threshold is an absolute difference:
$\text{RR}=0.986$ (0.956 to 1.017) \\
\bottomrule
\end{tabularx}}

\tpart{c}{Result}
{\small
\begin{tabularx}{\linewidth}{@{}K{3.4cm} Y@{}}
\toprule
Estimate & $-0.16$\pp\ (11.36\% motorcycle, 11.52\% car) \\
95\% confidence interval & $-0.52$ to $+0.19$\pp \\
Test statistic (d.f.) & $z=-0.89$ (normal) \\
p-value & 0.81 \\
Decision & fail to reject $H_0$ \\
\bottomrule
\end{tabularx}}

There is no evidence that motorcycle crashes are more severe here. The CI is narrow and lies far
below 5\pp, so this is a confident ``no meaningful difference'', not a lack of data. That
contradicts national statistics, in which motorcyclists are heavily over-represented in KSI.
Figure~\ref{fig:cramer} explains why: in this file \code{Vehicle\_Type} is almost unrelated to
every other column (Cram\'er's $V\le0.012$ with severity, area, light, road type and number of
vehicles), and the vehicle mix in single-vehicle crashes matches that in multi-vehicle crashes
to within 0.1\pp\ (e.g.\ buses 3.93\% vs 3.92\%). The column seems to record one vehicle per
accident, not the vehicle of the person who was hurt.

%------------------------------------------------------------ Test 4
\subsection*{Test 4 (PS2): do motorcycle crashes hurt fewer people?}

\tpart{a}{Hypotheses, written before running the test}
{\small
\begin{tabularx}{\linewidth}{@{}K{3.4cm} Y@{}}
\toprule
$H_0$ & $\mu_M-\mu_C\ge0$ \\
$H_1$ & $\mu_M-\mu_C<0$ \\
Significance level & $\alpha=0.05$ \\
Side & one-sided (left) \\
Test & Welch $t$-test (unequal variances) \\
Smallest difference that matters & 0.10 casualties per crash \\
\bottomrule
\end{tabularx}}

$\mu_M$, $\mu_C$: mean casualties per two-vehicle crash. 0.10 means 10 fewer injured people per
100 crashes, enough to change ambulance and hospital load.

\tpart{b}{Other tests we considered}
Two independent groups and a mean of counts. With $n$ in the tens of thousands, the CLT makes
the sample means normal despite the skew (Figure~\ref{fig:cas}), and $\sigma$ is unknown, so we
use a two-sample $t$-test. We chose Welch's version because the spreads differ.

{\footnotesize
\begin{tabularx}{\linewidth}{@{}L{2.4cm} Y@{}}
\toprule
\textbf{Also considered} & \textbf{Why we did not use it} \\
\midrule
Pooled $t$-test & assumes equal variances, rejected by Levene ($W=11.0$, $p=0.0009$) and by the
$F$-test ($F=0.93$, $p=10^{-17}$). Run anyway: $t=-3.32$, $p=0.0005$, same decision \\
Two-sample $z$-test & needs the population $\sigma$; at this $n$ it gives the same $z=-3.41$,
$p=0.0003$ \\
Mann--Whitney $U$ & 74\% of values are tied at 1, and it tests a shift in distribution, not the
mean. Robustness run: $p=0.001$, $P(M>C)+\tfrac12P(\text{tie})=0.496$, practically a coin
flip \\
Poisson rate test & casualties start at 1, and their variance is \emph{below} the mean
(var/mean $=0.44$ for $M$, $0.47$ for $C$); shifted to start at 0 they are overdispersed (1.65,
1.70). Either way the Poisson variance is wrong; it gives $p=0.012$ \\
\bottomrule
\end{tabularx}}

\tpart{c}{Result}
{\small
\begin{tabularx}{\linewidth}{@{}K{3.4cm} Y@{}}
\toprule
Estimate & $-0.015$ casualties per crash (1.362 motorcycle, 1.377 car) \\
95\% confidence interval & $-0.024$ to $-0.006$ \\
Test statistic (d.f.) & $t=-3.41$ (42{,}493) \\
p-value & $3.2\times10^{-4}$ \\
Decision & reject $H_0$ \\
\bottomrule
\end{tabularx}}

Motorcycle crashes do hurt fewer people on average, but \textbf{it does not matter}: the whole
CI lies inside $\pm0.10$, and the gap is 1.5 people per 100 crashes. With $n=341{,}739$, a
negligible gap is still significant. Dropping the 67 crashes with more than 10 casualties changes
nothing ($-0.015$, $p=3.3\times10^{-4}$).

\minihead{PS2 answer, and the multiple-testing check}
The expected tension between Tests 3 and 4 does not appear: motorcycle crashes are neither more
severe nor meaningfully smaller. Given Figure~\ref{fig:cramer}, we read this as a property of
how \code{Vehicle\_Type} was recorded, not as evidence that motorcycles are safe. Across the four
tests a Holm correction (adjusted $p$: $1.3\times10^{-15}$, $9.8\times10^{-8}$, $0.81$,
$6.5\times10^{-4}$ for Tests 1--4) \textbf{changes no decision}.

%-----------------------------------------------------------------------
\section{Limitations}

\begin{itemize}
  \item \textbf{No speed-limit column (the main confounder for both statements).} In Great
  Britain a road with street lighting defaults to 30\,mph, while an unlit rural single
  carriageway defaults to 60\,mph. Part of the 3.5\pp\ gap in PS1 may be speed, not light, and
  PS2 cannot be adjusted for speed either.
  \item \textbf{Association, not cause.} Lights are not placed at random: they go on junctions
  and built-up stretches. Only a before-and-after study of newly lit roads could show cause.
  \item \textbf{\code{Vehicle\_Type} looks detached from the crash.} Its near-zero association
  with every other column means PS2's conclusions describe this file, not motorcyclists in Great
  Britain.
  \item \textbf{Injury crashes only.} STATS19 records crashes with at least one injury that were
  reported to police, so the KSI share is a share of injury crashes, and slight injuries are
  under-reported.
  \item \textbf{Districts kept.} Test 2 uses 247 of 396 districts; the 149 with fewer than 20
  crashes in a group are left out, and the result may not apply to them.
  \item \textbf{Shifted dates.} Without usable dates we cannot check seasonal confounding (dark
  evenings in winter coincide with wet roads); the CMH check by road surface covers only part of
  this.
\end{itemize}

\end{document}
